% !TEX root = ../math601-notes.tex

\begin{defn*}
    For an $R$-module $M$, where $R$ is a commutative ring, we defined \textbf{tensor powers} $T^nM$ to be $M \otimes \overset{n}{\cdots} \otimes M$, and \textbf{symmetric powers} $S^nM$ to be $T^nM/I$. Moreover, we define \textbf{tensor algebras} to be $T^\bullet M = \bigoplus_{n \geq 0}T^nM$ and \textbf{symmetric algebras} to be $S^\bullet M = \bigoplus_{n \geq 0}S^nM$
\end{defn*}

\section{Exterior Powers and Exterior Algebras}

\begin{defn}
    For $n\geq 1$, we define the $n$-th \textbf{exterior power} $\bigwedge^nM$ to be an $R$-module, together with an $R$-multilinear map $\zeta: M \times \overset{n}{\cdots} \times M \to \Wedge^nM$ which is \textbf{alternating}. The map sends $(x_1,\dots,x_n) \mapsto x_1 \land \cdots \land x_n$.

    \textbf{Alternating} means $\zeta (x_1,\dots,x_n) = 0$ whenever $x_i = x_j$ for some $i \neq j$.
\end{defn}

\begin{prop}
    Alternating implies (but is not equivalent to) the condition that
    $$\zeta ( x_{\sigma(1)},\dots,x_{\sigma(n)})= \sgn (\sigma) \zeta(x_1,\dots,x_n)$$
    for all $\sigma \in S_n$
\end{prop}

\begin{prop}
    Let $n=2$. Given a bilinear map $\varphi :M \times M \to A$ for some $R$-module $A$, we have
    $$\varphi(x+y,x+y)=\varphi(x,x) + \varphi(x,y) + \varphi (y,x) + \varphi(y,y).$$
    If $\varphi (x,x)=0$ for all $x \in M$, then this implies $\varphi ( y,x) = - \varphi(x,y)$ for all $x,y \in R$.

    The converse is true if $2 \varphi(x,x) = 0 \implies \varphi (x,x) = 0$. This holds if multiplication by $2$ is injective in the $R$-module $A$.
\end{prop}

\begin{thm}
    The pair $(\Wedge^nM,\zeta)$ satisfies the universal property that for any $R$-module $A$ and alternating multilinear map $f:M \times \overset{n}{\cdots} \times M \to A$, there exists a unique $R$-linear map $\tilde f : \Wedge^nM \to A$ such that $\tilde f \circ \zeta = f$. That is, the diagram below commutes:

    % https://q.uiver.app/#q=WzAsMyxbMCwwLCJNIFxcdGltZXNcXG92ZXJzZXR7bn17XFxjZG90c31cXHRpbWVzIE0iXSxbMiwwLCJcXGJpZ3dlZGdlIF5uTSJdLFsyLDIsIkEiXSxbMCwxLCJcXHpldGEiXSxbMSwyLCJcXGV4aXN0cyAhIFxcdGlsZGUgZiIsMCx7InN0eWxlIjp7ImJvZHkiOnsibmFtZSI6ImRhc2hlZCJ9fX1dLFswLDIsImYiLDJdXQ==
    \[\begin{tikzcd}[ampersand replacement=\&]
    	{M \times\overset{n}{\cdots}\times M} \&\& {\bigwedge ^nM} \\
    	\\
    	\&\& A
    	\arrow["\zeta", from=1-1, to=1-3]
    	\arrow["f"', from=1-1, to=3-3]
    	\arrow["{\exists ! \tilde f}", dashed, from=1-3, to=3-3]
    \end{tikzcd}\]
\end{thm}

\begin{proof}
    We start with existence. Let $\Wedge ^nM = T^nM/J$, where $J$ is the $R$-submodule of $T^nM = M \otimes \overset{n}{\cdots} \otimes M$ generated by the tensors $x_1 \otimes \cdots \otimes x_n$, where $x_i = x_j$ for some $i \neq j$. \textbf{Exercise:} check this.
\end{proof}

\begin{rmk}
    We have $\Wedge ^1 M = M$, and $\Wedge^0M = R$ by convention.
\end{rmk}

Similar to the case of symmetric powers, we have bilinear maps $\Wedge^aM \otimes \Wedge^bM  \to \Wedge^{a+b}M$, where
$$(x_1 \land \cdots \land x_a) \otimes (y_1 \wedge \cdots \wedge y_b) \mapsto x_1 \land \cdots \land x_a \land y_1 \land \cdots \land y_b$$
\begin{defn}
    These make the $R$-module
    $$\Wedge ^\bullet M := \bigoplus _{n=0}^\infty \Wedge^nM$$
    into an $R$-algebra called the \textbf{exterior algebra} of $M$. $\Wedge^\bullet M$ is a graded, skew-commutative $R$-algebra.

    \textbf{Skew-commutative} means that for $x \in \Wedge^a M,y \in \Wedge ^bM$, we have
    $$y \land x = (-1)^{ab}x \land y$$
    (To see this, think about how $x$ has $a$ wedges, $y$ has $b$ wedges, and the sign flips every time they cross.)
\end{defn}

\begin{thm}
    For any $R$-algebra $A$ and map $f : M \to A$ of $R$-modules such that $f(x)^2=0$ in $A$ for any $x \in M$, there exists a unique homomorphism of $R$-algebras $\tilde f : \Wedge ^\bullet M \to A$, making the below diagram commute:

    % https://q.uiver.app/#q=WzAsMyxbMCwwLCJNIl0sWzIsMCwiXFxMYW1iZGEgXlxcY2lyYyBNIl0sWzEsMiwiQSJdLFswLDIsImYiLDJdLFsxLDIsIlxcdGlsZGUgZiIsMCx7InN0eWxlIjp7ImJvZHkiOnsibmFtZSI6ImRhc2hlZCJ9fX1dLFswLDFdXQ==
    \begin{tikzcd}[ampersand replacement=\&]
    	M \&\& {\Wedge ^\bullet M} \\
    	\\
    	\& A
    	\arrow[from=1-1, to=1-3]
    	\arrow["f"', from=1-1, to=3-2]
    	\arrow["{\tilde f}", dashed, from=1-3, to=3-2]
    \end{tikzcd}
\end{thm}

\begin{proof}
    Construction: Let
    $$\Wedge ^\bullet M := T^\bullet M / \angles{ 2\text{-sided ideals generated by }x \otimes x \in T^2M, \forall x \in M }$$
\end{proof}

\begin{prop}
    $$\Wedge^n(M \oplus N) \cong \bigoplus _{a+b=n}\Wedge^aM \otimes \Wedge ^bN$$
\end{prop}

\begin{proof}
    We have maps $\Wedge ^aM \otimes \Wedge^bN \to \Wedge ^{a+b}(M \oplus N)$ which map
    $$(x_1 \land \cdots \land x_a) \otimes ( y_1 \land \cdots \land y_b) \mapsto (x_1 \land \cdots \land x_a) \land (y_1 \land \cdots \land y_b).$$
    Here, $x \in M$ is defined with $(x,0) \in M \oplus N$, and $y \in N \equiv (0,y) \in M \oplus N$. That is, we have a natural inclusion $M \hookrightarrow M \oplus N, N \hookrightarrow M \oplus N$.

    Adding these together produces a map
    $$g : \bigoplus _{a+b = n}\Wedge^aM \otimes \Wedge ^bN \to \Wedge ^n(M \oplus N).$$
    The inverse of $g$ is harder to define. For $a,b \geq 0$ such that $a+b=n$, we will consider partitions of the set $[1,n] := \{ 1,2,\dots,n\}$ into (disjoint) subsets $I$ and $J$, with $|I| = a, |J|=b$. We write
    $$I = \{ i_1<\cdots <i_a\},J=\{j_1 < \cdots < j_b\},I \sqcup J = [1,n].$$
    For each pair $(I,J)$, let $sgn(I,J)$ be the sign of the permutation
    $$\matr{1 & 2 & \cdots & n \\ i _1 & i_2 \cdots i_a & \cdots & j_1 \cdots j_b} \in S_n.$$
    Consider the $R$-multilinear map
    $$f: ( M \oplus N) \times \cdots \times (M \oplus N) \to \bigoplus_{a+b = n}\Wedge^aM \otimes \Wedge^bN$$
    given by
    $$f((x_1,y_1),(x_2,y_2),\dots,(x_n,y_n))$$
    $$=\bigoplus_{a+b=n}\paren{\sum_{I \sqcup J = [1,n],|I|=a,|J|=b }\sgn(I,J)(x_{i_1}\land \cdots \land x_{i_a})\otimes (y_{j_1} \land \cdots \land y_{j_b}}$$
    For example, when $n=2$, we have
    $$f((x_1,y_1),(x_2,y_2)) = x_1 \land x_2 + x_1 \otimes y_2 - x_2 \otimes y_1 + y_1 \land y_2.$$
    To understand the sign, consider the identity:
    $$(x_1 + y_1)\land (x_2+y_2) = x_1\land x_2 + x_1 \land y_2 + y_1 \land x_2 + y_1 \land y_2.$$
\end{proof}

\begin{ex}
    When $n=2$, we have
    $$\Wedge ^2(M \oplus N) \cong (\Wedge ^2 M) \oplus (M \otimes N) \oplus (\Wedge^2N)$$
\end{ex}

\begin{ex}
    For $M =R$, we have $\Wedge ^nR=0$ if $n >1$. This is because if $a,b \in R$, we have $a \land b = ab(1 \land 1) =0$.
\end{ex}

\begin{cor}
    If $M$ is a free $R$-module with basis $e_1,\dots,e_n$, then $\Wedge^kM$ is free with basis $e_{i_1}\land \cdots \land e_{i_k}$ for $1 \leq i_1 < \cdots < i_k \leq n$. So, $\Wedge ^kM$ has rank $\binom{n}{k}$.
\end{cor}
\begin{proof}
    We have $\Wedge ^0R = \Wedge ^1R=R$, and $\Wedge^ aR =0$ if $a>1$. Since $M$ is free on $e_1,\dots,e_n$, we have
    $$ M \cong \bigoplus_{i=1}^nRe_i.$$
    The previous proposition and induction give
    $$\Wedge^kM = \Wedge^k(Re_1\oplus \cdots \oplus Re_n) \cong \bigoplus _{a_1+\cdots + a_n=k,a_i \geq 0}\Wedge^{a_1}Re_1 \otimes \cdots \otimes \Wedge^{a_n}Re_n$$
    $$\cong \bigoplus _{a_1+\cdots + a_n =k,a_i \in \{0,1\}} \Wedge^{a_1}Re_1 \otimes \cdots \otimes \Wedge^{a_n}Re_n$$
    as $R$-modules. This is free of rank $\binom{n}{k}$, with basis $1 \otimes \cdots \otimes e_{i_1}\otimes \cdots \otimes 1 \otimes e_{i_k}\otimes 1 \otimes \cdots \otimes 1$, and $\{ i_1,\dots,i_k\} \subset [1,n]$. The inverse isomorphism maps these to the elements $e_{i_1}\land \cdots \land e_{i_k}$ in $\Wedge^kM$.
\end{proof}

\begin{rmk}
    The definition of tensor, symmetric, and exterior powers are \textbf{functorial}: given a map $f : M \to N$ of $R$-modules, we get maps $T^nf:T^nM \to T^nN$, $S^nf : S^nM \to S^nN$ and $\Wedge ^nf : \Wedge ^nM \to \Wedge ^nN$ for all $n \geq 0$.

    For example, $\Wedge^nf(x_1\land \cdots \land x_n) := f(x_1) \land \cdots \land f(x_n)$
\end{rmk}

% https://q.uiver.app/#q=WzAsNSxbMCwwLCIoeF8xLFxcZG90cyx4X24pIl0sWzMsMCwiZih4XzEpIFxcbGFuZCBcXGNkb3RzIFxcbGFuZCBmKHhfbikiXSxbMCwyLCJUXm5NIl0sWzAsMywieF8xIFxcb3RpbWVzIFxcY2RvdHMgXFxvdGltZXMgeF9uIl0sWzAsNCwiXFxMYW1iZGEgXm5mIl0sWzAsMiwiXFx0aWxkZSBmXnsobil9IiwyXSxbMyw0XSxbMCwxLCIiLDAseyJzdHlsZSI6eyJ0YWlsIjp7Im5hbWUiOiJtYXBzIHRvIn19fV0sWzIsMV0sWzQsMV1d
\begin{tikzcd}
	{(x_1,\dots,x_n)} &&& {f(x_1) \land \cdots \land f(x_n)} \\
	\\
	{T^nM} \\
	{x_1 \otimes \cdots \otimes x_n} \\
	{\Wedge ^nf}
	\arrow[maps to, from=1-1, to=1-4]
	\arrow["{\tilde f^{(n)}}"', from=1-1, to=3-1]
	\arrow[from=3-1, to=1-4]
	\arrow[from=4-1, to=5-1]
	\arrow[from=5-1, to=1-4]
\end{tikzcd}

(Note: I'm not sure if this is correct)

\begin{cor}
    If the $R$-module $M$ has $n$ generators, then $\Wedge ^kM = 0$ for all $k > n$.
\end{cor}

\begin{proof}
    To say that $M$ has $n$ generators is equivalent to saying that there is a surjective $R$-linear map $\pi: F \to M$, where
    $$F = \bigoplus_{i=1}^nRe_i \cong R^n$$
    is a free $R$-module of rank $n$. We get an induced map $\Wedge ^k \pi : \Wedge ^kF \to \Wedge ^kM$ which is also surjective [In fact, if $f: M \to N$ is surjective, then $T^nf,S^nf,\Wedge^nf$ are all surjective]. If $k>n$, then the corollary gives $\Wedge ^kF = 0$, hence $\Wedge ^kM =0$ also.
\end{proof}

\begin{notebox*}[Challenge]
    If $M$ and $N$ are free $R$-modules, and $f:M \hookrightarrow N$ is injective, what about $T^nf,S^nf,$ and $\Wedge^nf$?
\end{notebox*}

\begin{ex}
    For $M$ a free $R$-module of rank $n$ with basis $e_1,\dots,e_n$, $\Wedge^\bullet M$ is free of rank $2^n$ with basis $e_{i_1}\land \cdots \land e_{i_k}$ for $1 \leq i_1 < \cdots < i_k \leq n$. As an algebra, $\Wedge ^\bullet M$ is generated by $e_1,\dots ,e_n$ modulo the relations $e_i \land e_i=0$ and $e_i \land e_j + e_j \land e_i =0$ for $i,j \in [1,n]$. (Recall: $S^\bullet M \equiv R[e_1,\dots,e_n], T^\bullet M \equiv R \angles {e_1,\dots,e_n}$).
\end{ex}

\section{Determinants}

\begin{defn}
    Let $F$ be a free $R$-module on $n$ generators, and $T:F \to F$ be an $R$-linear map. Then, we get an induced map $\Wedge^nT: \Wedge^nF \to \Wedge^nF$. Here, $\Wedge^nF$ is free on one generator (namely $e_1 \land \cdots \land e_n$, where $e_1,\dots,e_n$ is an $R$-basis of $F)$. This implies that $\Wedge^nT(x) = \lambda x$ for all $x \in \Wedge^nF$, and for some $\lambda \in R$. This $\lambda$ is called the \textbf{determinant} of $T$, denoted $\det T$.
\end{defn}

\begin{notebox*}
    The determinant $\det T$ is defined independently of the choice of basis. So, we have
    $$\Wedge^nT(x) = (\det T)\cdot x$$
    for all $x \in \Wedge^nF$.
\end{notebox*}

\begin{ex}[Formula for $\det T$]
    Let $e_1,\dots,e_n$ be a basis for $F$. For any $\sigma \in S_n$, we can think of $\sigma$ as a linear map $\sigma:F \to F$ given by $\sigma (e_i):=e_{\sigma (i)}$. Then, we have
    $$\Wedge^n\sigma (e_1\land \cdots \land e_n) = \sigma (e_1) \land \cdots \land \sigma (e_n)= \sgn (\sigma)\cdot e_1\land \cdots \land e_n,$$
    hence $\det (\sigma ) = \sgn (\sigma)$.
\end{ex}

\begin{prop}[Formula for $\det T$]
    In general, if $T:F \to F$ is an $R$-linear map, write $T(e_j) = \sum_{i=1}^na_{ij}e_i$, $1 \leq j \leq n$. The matrix $A:= \{a_{ij}\}_{1 \leq i,j \leq n}$ is $(T)_E$, i.e. the matrix representing $T$ with respect to the ordered basis $E:=(e_1,\dots,e_n)$. Then,
    $$\det (T) = \sum_{\sigma \in S_n}\sgn(\sigma) a_{\sigma(1), 1}a_{\sigma(2), 2}\cdots a_{\sigma(n), n}$$
\end{prop}

\begin{proof}
    We have
    $$\Wedge^nT(e_1 \land \cdots \land e_n) = T(e_1) \land \cdots \land T(e_n) = \sum_{i_1}a_{i_11}e_{i_1}\land \sum_{i_2}a_{i_22}e_{i_2}\land \cdots \land \sum_{i_n}a_{i_nn}e_{i_n}$$
    $$= \sum_{i_1,\dots,i_n}a_{i_11}a_{i_22}\cdots a_{i_nn}\cdot e_{i_1}\land \cdots \land e_{i_n} = \sum_{\sigma \in S_n}a_{\sigma (1)1}a_{\sigma (2)2}\cdots a_{\sigma (n)n}\cdot e_{\sigma (1)}\land \cdots \land e_{\sigma (n)}$$
    $$= \paren{ \sum_{\sigma \in S_n}\sgn (n) a_{\sigma (1)1}\cdots a_{\sigma (n)n}}e_1 \land \cdots \land e_n$$
\end{proof}

\begin{rmk}
    We denote $\det T$ by $\det A$ for any $A \in M_n(R)$ representing $T$.
\end{rmk}

\begin{ex}
    $$\det \matr{a_{11} & a_{12} \\ a_{21} & a_{22}}=a_{11}a_{22}-a_{21}a_{12}$$
    where the first term corresponds to the identity permutation, and the second corresponds to the permutation $(1\:\:2) \in S_2$.
\end{ex}

\begin{prop}
    If $S,T:F \to F$ are 2 $R$-linear maps, then
    $$\det (T \circ S) = \det (T) \cdot \det (S)$$
\end{prop}
\begin{proof}
    % https://q.uiver.app/#q=WzAsNSxbMCwwLCJcXExhbWJkYV5uRiJdLFsyLDAsIlxcTGFtYmRhXm5GIl0sWzQsMCwiXFxMYW1iZGFebkYiXSxbMywxXSxbMSwxXSxbMCwxLCJcXExhbWJkYV5uUyJdLFsxLDIsIlxcTGFtYmRhXm5UIl0sWzMsMl0sWzAsNF0sWzQsMywiXFxMYW1iZGFebihUXFxjaXJjIFMpIiwyXV0=
    \begin{tikzcd}[ampersand replacement=\&]
    	{\Wedge^nF} \&\& {\Wedge^nF} \&\& {\Wedge^nF} \\
    	\& {} \&\& {}
    	\arrow["{\Wedge^nS}", from=1-1, to=1-3]
    	\arrow[from=1-1, to=2-2]
    	\arrow["{\Wedge^nT}", from=1-3, to=1-5]
    	\arrow["{\Wedge^n(T\circ S)}"', from=2-2, to=2-4]
    	\arrow[from=2-4, to=1-5]
    \end{tikzcd}
    commutes.
\end{proof}
\begin{prop}
    For any $A \in M_n(R)$, we have $\det (A^t) = \det (A)$.
\end{prop}

\begin{proof}
    If $A=\{a_{ij}\}$, then $A^t=\{b_{ij}\}$ where $b_{ij}=a_{ji}$. Therefore,
    $$\det (A^t) = \sum_{\sigma \in S_n}\sgn (\sigma) a_{1 \sigma (1)}a_{2 \sigma (2)}\cdots a_{n \sigma (n)} = \sum_{\sigma \in S_n}\sgn (\sigma) a_{\sigma ^{-1}(1) 1}\cdots a_{\sigma ^{-1}(n)n}$$
    $$\sum_{\sigma \in S_n}\sgn (\sigma^{-1}) a_{\sigma ^{-1}(1) 1}\cdots a_{\sigma ^{-1}(n)n}=\det (A)$$
\end{proof}

\begin{thm}[Cramer's Rule]
    Let $A = \{a_{ij}\} \in M_{n}(R)$. If $A \cdot \matr{x_1 \\ \vdots \\ x_n}= \matr{b_1 \\ \vdots \\ b_n}$ with $a_{ij},x_i,b_i \in R$, then we have
    $$\det (A) \cdot x_i = \det \matr{a_{11} & \cdots & b_1 & \cdots & a_{1n} \\ \vdots & & \vdots && \vdots \\ a_{1n}& \cdots & b_n & \cdots & a_{nn}}$$
    where the $b$ column is the $i$-th column.
\end{thm}

\begin{proof}
    Let $F= R^n$, and regard its elements as column vectors. Let $A = (\vec a_1 \cdots \vec a_n)$ where $\vec a_1,\dots, \vec a_n$ are the columns of $A$, and $\vec x,\vec b=\matr{b_1 \\ \vdots \\ b_n} \in F$. Then, $A \vec x = \vec b$ implies that $\vec b = x_1 \vec a_1 \cdots x_n \vec a_n$. Therefore,
    $$\det (\vec a_1 \cdots \vec b \cdots \vec a_n) = \det ( \vec a_1 \cdots \sum_{j=1}^nx_j\vec a_j \cdots \vec a_n)=\sum_{j=1}^nx_j \det (\vec a_1 \cdots \vec a_j \cdots \vec a_n) = \det (A) x_i.$$
\end{proof}

\begin{cor}
    If $\det T$ is not a zero divisor in $R$, then $T$ is injective.
\end{cor}

\begin{proof}
    By Cramer's Rule,
    $$x = \sum_{i=1}^nx_ie_i \in \ker T\iff T(x) =0$$
    implies that $\det (T) \cdot x_i =0$ for all $1 \leq i \leq n$. Since $\det (T)$ is not a zero divisor in $R$, we get $x_i =0, 1\leq i \leq n$.
\end{proof}

\begin{notebox*}[Exercises]
    \begin{enumerate}
        \item Prove the converse to the above corollary.
        \item Prove the usual Laplace expansion of $\det (A)$ along rows and columns. A generalization is on the homework.
    \end{enumerate}
\end{notebox*}

\begin{rmk}
    \begin{enumerate}
        \item Here is a different method of defining determinants:
    \end{enumerate}

    Let $M := R^n$. Then, $\det :M \times \overset{n}{\cdots} \times M \to R$ is the unique alternating multilinear map such that $\det (I_n)=1$. This defines $\det (A)$ for any $A \in M_n(R)$
    \begin{enumerate}
        \item[2.] Determinants arise naturally as signed volumes:
    \end{enumerate}
    For $v_1,\dots,v_n \in \R^n$, let $V_S(v_1,\dots,v_n)$ denote the signed volume of the polytope
    $$P:= \brax{\sum_{i=1}^n \lambda _i v_i \mid 0 \leq \lambda _i \leq 1}$$
    with "sides" $v_1,\dots,v_n$. Precisely, we have
    $$V_S(v_1,\dots,v_n) = \case{0 & \text{if } v_1,\dots,v_n \text{ are linearly dependent} \\ Vol(P) & \text{if } v_1,\dots,v_n \text{ are a positively oriented basis }\\ - Vol(P) & \text{otherwise}}$$
    where positively oriented means $\det (v_1\cdots v_n)>0$. Then,
    $$Vol(P) = |\det (\vec v_1 \cdots \vec v_n)|$$
\end{rmk}

\section{Localization of Rings and Modules}

\begin{defn}
    Let $R$ be a commutative ring. A subset $S$ of $R$ with $1 \in S$ is \textbf{multiplicatively closed} if $x,y \in S \implies xy \in S$. For such $S$, there exists a commutative ring $S^{-1}R$ with a map of rings $\pi : R \to S^{-1}R$ such that the following universal property holds:

    For any ring homomorphism $f : R \to R'$ such that $f(s)$ is a unit in $R'$ for each $s \in S$, there exists a unique map $\tilde f:S^{-1}R \to R'$ such that $\tilde f \circ \pi = f$. That is, the below diagram commutes:

    % https://q.uiver.app/#q=WzAsMyxbMCwwLCJSIl0sWzIsMCwiU157LTF9UiJdLFsxLDEsIlInIl0sWzAsMSwiXFxwaSJdLFswLDIsImYiLDJdLFsxLDIsIlxcZXhpc3RzICEgXFx0aWxkZSBmIiwwLHsic3R5bGUiOnsiYm9keSI6eyJuYW1lIjoiZGFzaGVkIn19fV1d
    \begin{tikzcd}[ampersand replacement=\&]
    	R \&\& {S^{-1}R} \\
    	\& {R'}
    	\arrow["\pi", from=1-1, to=1-3]
    	\arrow["f"', from=1-1, to=2-2]
    	\arrow["{\exists ! \tilde f}", dashed, from=1-3, to=2-2]
    \end{tikzcd}
\end{defn}

\begin{proof}
    Construction:
    $$S^{-1}R:= R \times S/ \sim$$
    where $(r,s) \sim (r',s') \iff u(rs'-sr')=0$ for some $u \in S$. Denote the equivalence classes of $(r,s)$ by $\frac{r}{s}$, and define the operations in $S^{-1}R$ by the usual addition and multiplication rules of fractions. Then, the map $\pi :R \to S^{-1}R$ sends $r$ to $\frac{r}{1}$.

    \begin{rmk}
        $$\ker \pi = \{ r \in R \mid sr =0\text{ for some } s \in S\}$$
        so $\pi$ is injective
    \end{rmk}
    $S$ contains neither zero nor any zero divisor in $R$. Also, $S^{-1}R=0$ if $0 \in S$.
\end{proof}

\begin{defn}
    This definition extends to modules $M$ over $R$ (with $S$ still a subring of $R$). We define $S^{-1}M$ as a set of equivalence classes $[\frac{m}{s}]$, where $m \in M$ and $s \in S$, and where $\frac{m}{s} \sim \frac{m'}{s'} \iff s'' (s'm-sm')=0$ for some $s'' \in S$.
\end{defn}

\begin{rmk}
    There is a map $\pi :M \to S^{-1}M$ with $\pi (m) = [\frac{m}{1}]$ for all $m \in M$. Also,
    $$\ker \pi = \{ m \in M \mid sm =0 \text{ for some } s \in S\}.$$
\end{rmk}

\begin{thm}
    If $M$ and $N$ are $R$-modules, and $f:M \to N$ is a module homomorphism, then for $S \subset R$ multiplicatively closed, we get an induced $S^{-1}R$-module map $S^{-1}M \to S^{-1}N$ sending $\frac{m}{s} \mapsto \frac{f(m)}{s}$.
\end{thm}

\begin{prop}
    $$S^{-1}M \cong S^{-1}R \otimes _RM$$
    as $S^{-1}R$ modules. That is, $S^{-1}M$ is the $S^{-1}R$ module obtained by extension of scalars from the $R$-module $M$.
\end{prop}

\begin{proof}
    We have a map $S^{-1}R \otimes M \to S^{-1} M$ sending $\frac{r}{s}\otimes m \mapsto \frac{rm}{s}$ with inverse $S^{-1}M \to S^{-1}R \otimes M$ sending $\frac{m}{s}\mapsto \frac{1}{s}\otimes m$. To see that the latter map is well defined, note that if $\frac{m}{s}=\frac{m'}{s'}$ in $S^{-1}M$, then $x(s'm - sm')=0$ for some $x \in S$. So,
    $$\frac{1}{s}\otimes m = \frac{1}{xs's} \otimes (s'xm)=\frac{1}{xs's}\otimes (xsm')=\frac{1}{s}\otimes m'$$
\end{proof}

We'll use these constructions when $S^{-1}R$ is the fraction field $K$ of an integral domain (so when $S=R\setminus \{0\}$).

\begin{notebox*}[Exercise]
    Show that
    $$S^{-1}(M \oplus N) \cong S^{-1}M \oplus S^{-1}N$$
    for any two $R$-modules $M$ and $N$.
\end{notebox*}

\section{Finitely Generated Modules Over a PID}

\begin{notebox*}[Goal]
    Classify all finitely generated modules over a PID, up to isomorphism.
\end{notebox*}

\begin{ex}
    When $R=\Z$, this gives a classification of all finitely generated abelian groups.
\end{ex}

\begin{defn}
    Given any $R$-module $M$, the \textbf{torsion submodule} of $M$, denoted $M_{\tors}$, is defined by
    $$M_{\tors}=\{x \in M \mid rx =0 \text{ for some non-zero } r \in R\}$$
\end{defn}

\begin{defn}
    We say that $M$ is \textbf{torsion free} if $M_{\tors} = 0$.
\end{defn}

\begin{ex}
    \begin{enumerate}
        \item Any free $R$-module $R^n$ is torsion free, provided $R$ is a domain.
        \item The converse of (1) is false, e.g. $\Q$ is a torsion free $\Z$-module, but $\Q$ is not free as a $\Z$-module.
        \item If $M = \Z^{10}\oplus \Z/2\Z \oplus \Z/8\Z$, then $M_{\tors}=\Z/2\Z \oplus \Z/8\Z$.
    \end{enumerate}
\end{ex}

\begin{prop}
    We have an exact sequence
    $$0 \to M_{\tors} \to M \to M/M_{\tors} \to 0$$
    for any module $M$. The quotient $M/M_{\tors}$ is torsion free.
\end{prop}

Suppose that $R$ is an integral domain with quotient field $K$. That is, let $K = S^{-1}R$, where $S = R \setminus \{0\}$. Moreover, for any $R$-module $M$, we set
$$M_K := S^{-1}M,$$
which is a $K$-vector space. The natural map $\alpha : M \to M_K$ defined by $m \mapsto \braks{\frac{m}{1}}$ has
$$\ker \alpha = \{ m \in M \mid \frac{m}{1} \sim \frac{0}{1}\} = \{ m \in M \mid \exists r \in R \setminus\{0\}, rm=0\} = M_{\tors}.$$
Hence, $M$ is torsion free $\iff$ $M \overset{\alpha}{\hookrightarrow} M_K$ as $R$-modules.

\begin{ex}
    \begin{enumerate}
        \item If $M=R,$ then $M_K=K$. More generally,, if $M = R^n$ is free of rank $n$, then $M_K = K^n$
        \item Suppose $M = \Z/n\Z$ as a $\Z$-module. Then, $M_\Q = \{0 \}$. Indeed, for each $k \in \Z^\times,$ we have $\frac{x}{k}=\frac{nx}{nk}=\frac{0}{nk}=0$ for all $x \in M$. More generally, if $M = \Z^n \oplus \Z/k_z\Z \oplus \cdots \oplus \Z/k_r\Z,$ then $M_\Q \cong \Q^n$.
    \end{enumerate}
\end{ex}
\begin{defn}
    Let $R$ be a domain, and $M$ and $R$-module. The \textbf{rank} of $M$, denote $\rank(M)$, is the dimension of $M_K$ as a $K$-vector space.
\end{defn}

\begin{rmk}
    If $M$ is finitely generated, then $\rank(M)$ is finite.
\end{rmk}

\begin{thm}
    Suppose $R$ is a PID, and $M$ is a finitely generated torsion-free $R$-module. Then, $M$ is free.
\end{thm}

\begin{proof}
    We induct on $\rank(M)$. If $\rank(M)=0$, then $M=0$, since $M$ is torsion free. Now, suppose $\rank(M) >0$. Then, $M_K$ is a finite dimensional non-zero $K$-vector space, so there exists a surjective map $M_K \to K$. We then have $M \overset{\alpha}{\hookrightarrow} M_K \to K$, and $\Img (\alpha)$ generates $M_K$. It follows that there exists a non-zero map $\beta :M \to K$. Let $M = \angles{m_1,\dots,m_r}$, so then multiplying $\beta$ by a suitable element in $R$ to clear denominators, we obtain a non-zero map $\varphi : M \to R$. Now, $\Img \varphi = \angles{b}$ is an ideal of $R$ (which is the same as a submodule of $R$). Choose $a \in M$ with $\varphi (a) = b$.

    Now, we claim $M \cong \ker \varphi \oplus Ra \cong \ker \varphi \oplus R$ as $R$-modules: Let $N:= \ker \varphi$. Then, given $m \in M$, we write $\varphi (m)=xb=x\varphi(a) = \varphi(xa)$, so $m - xa \in M$. Therefore, $m \in N + Ra$, so $M = N +Ra$. If $ra \in N$ for some $r \in R$, then $0=\varphi (ra) =r\varphi (a) = rb$, hence $r=0$, so $ra = 0$. Hence $N \cap Ra = \{0\}$, proving the claim.

    The claim implies that $M \cong N \oplus R$. Hence, $M_K \cong N_K \oplus K$, while $\rank (N_K) < \rank (M)$. We hence conclude by induction.
\end{proof}

\begin{prop}
    If $F$ is finitely generated and free over a PID $R$, and $N \subset F$ is a submodule. Then, there exists a basis $e_1,\dots,e_n$ of $F$, and $d_1\mid d_2 \mid \cdots \mid d_k \in R$ for $k \leq n$ (so $\angles{d_1}\supset \cdots \supset \angles{d_k}$) such that $N$ has a basis $d_1e_1,\dots,d_ke_k$. In particular, $N$ is free. Moreover, the integer $k$, and the ideals $\angles {d_1}\supset \cdots \supset \angles{d_k}$ are uniquely determined by $F$ and $N$.
\end{prop}

\begin{proof}
    We may assume $N \neq 0$. We use induction on $\rank (F)$. If $\rank (F) =1$, then $F=R$, so $N = (d) = Rd$ for some $d \in R$.

    In general, suppose $N \subset F$ is any submodule. The idea is that a general element of $N$ is $\sum_{i=1}^k r_id_i e_i=d_1y$ for some $y$. Hence, for all $\varphi \in F^\star := \Hom (F,R)$, $\varphi (N) \subset (d_1)$. Moreover, there exists $e_1^\star \in F^*$ such that $e_1^*(N) = (d_1)$. So, consider the set
    $$S := \{ \varphi(N)\mid \varphi \in F^*\}$$
    of ideals in $R$. Since $R$ is Noetherian (it's a PID), we know $S$ has a maximal element $=:I$. Write $I=\angles{d_1}$, and let $x \in N$ and $\varphi \in F^*$ be such that $\varphi (x) = d_1$.

    Now, let $x = \alpha _1e_1 + \cdots + \alpha _n e_n$. Then, consider $\angles{d_1,\alpha_i}= \angles {sd_1 + t \alpha _i}$ for some $s, t \in R$. Then,
    $$(s \varphi + te_i^*)(x) = sd_1 + t\alpha _i,$$
    so $(d_1) \subset (d_1,\alpha_i) \subset (s \varphi + te_i^*)(N)$. Now, since $\angles {d_1}$ is maximal in $S$, we deduce that $\angles {d_1}=\angles {d_1,\alpha _i}$, hence $d \mid \alpha _i$. Therefore, $x = d_1y$ for some $y \in F$. Note that $\varphi(y)=1$. Let $F' := \ker \varphi$. Arguing as in the previous proposition, we get $F \cong F' \oplus Ry$. Moreover, $N = N' \oplus Rx=N' \oplus Rd_1y$, where $N'= N \cap F'$.

    Since $F'$ is torsion free, we know by the previous proposition that $F'$ is free, and $\rank (F') = n-1$, while $N' \subset F'$. By induction on $\rank (F)$, we have a basis $e_2,\dots,e_n$ of $F'$, and $d_2 \mid d_3 \mid \cdots \mid d_k$ such that $N' = Rd_2e_2 \oplus \cdots \oplus Rd_ke_k$. Now, set $e_1:=y$. Then, $N = R d_1e_1 \oplus \cdots \oplus Rd_ke_k$. Moreover, we have $\angles {d_1} \supset \angles {d_2}$, because if not, we can produce $\psi \in F^*$ with $\psi (N) \supset \angles{d_1,d_2} \supsetneq \angles {d_1}$ as before, which gives a contradiction.  We have thus proven existence.

    For uniqueness, consider
    $$M:= F/N \cong R^{n-k} \oplus R/\angles{d_1}\oplus \cdots \oplus R/\angles{d_k}.$$
    Then, $n-k = \rank (M) = \rank (F/N)$, so $k$ is uniquely determined. Moreover,
    $$M_{\tors} = R/\angles {d_1}\oplus \cdots \oplus R/\angles{d_k}.$$

    \begin{defn}
        Given any $R$-module $T$, the \textbf{annihilator} of $T$ is the ideal of $R$ given by
        $$\Ann (T) := \{ r \in R \mid rt = 0\:\: \forall t \in T \}.$$
    \end{defn}
    Now,
    $$\angles {d_k} = \Ann (M_{\tors} ), \angles{d_{k-1}}= \Ann ( \Wedge ^2 M_{\tors}),\dots,\angles {d_{k+1-m}}= \Ann ( \Wedge^mM_{\tors}),\dots$$
    To see this, we use
    $$\Wedge^2(R/\angles d) = 0 \impliedby (r_1 + \angles d) \land (r_2 + \angles d) = r_1r_2(1+\angles d) \land (1+ \angles d)=0.$$
    Now, recall from last year that
    $$R/I \otimes _R R/J \cong R/(I+J)$$
    by the map $(r_1+I) \otimes (r_2 + I) \mapsto r_1r_2 + (I+J)$. So,
    \begin{align*}
    \Wedge^m M_{\tors} &= \Wedge^m \paren{ \bigoplus_{j=1}^k R/\angles {d_j}}=\bigoplus_{j_1<\cdots < j_m}R/\angles{d_{j_1}}\otimes \cdots \otimes R/\angles {d_{j_m}}
    \\ &= \bigoplus _{j_1< \cdots < j_m}R/[\angles {d_{j_1}}+ \cdots + \angles {d_{j_m}}] = \bigoplus_{j_1<\cdots < j_m}R/\angles {d_{j_1}}
    \\ & \implies \Ann ( \Wedge^m M_{\tors} ) = \angles {d_{k+1-m}}
    \end{align*}
\end{proof}

\begin{thm}[Structure Theorem (Invariant Factor Form)]
    Let $R$ be a PID, and $M$ a finitely generated $R$-module. Then, we have
    $$M \cong R^m \oplus R/\angles {d_1}\oplus \cdots \oplus R/\angles {d_k},$$
    where $\angles {d_1}\supset \cdots \supset \angles {d_k} \neq 0$. Moreover, $m,k$, and the ideals $\angles {d_1},\dots, \angles{d_k}$ are uniquely determined by $M$.
\end{thm}

\begin{proof}
    We can produce a surjection $\varphi: F \to M$, where $F$ is a free module. Let $N := \ker \varphi$, and apply the previous theorem to prove existence. To prove uniqueness, we see that $m = \rank (M)$ is uniquely determined. The rest follows as in the last proposition.
\end{proof}

\begin{defn}
    The integer $m$ is the \textbf{rank} of $M$ (and equals $\dim _K(M_K)$). The elements $d_1,\dots,d_k$ of $R$, defined up to units, are called the \textbf{invariant factors} of $M$.
\end{defn}

\begin{rmk}
    If
    $$d = \prod _{i=1}^r p_i^{k_i}$$
    is the factorization of $d \in R$ into primes, then by the Chinese Remainder Theorem
    $$R/\angles{d}\cong R/\angles{p_1^{k_1}}\times \cdots \times R/\angles{p_r^{k_r}}$$
    as rings. Using this in the last result gives the following theorem.
\end{rmk}

\begin{thm}[Structure Theorem (Elementary Divisor Form)]
    Given $M$, a finitely generated module over a PID $R$, we have
    $$M \cong R^m \oplus \bigoplus_{p\text{ prime}} \paren{ R/\angles{p^{n_1(p)}}\oplus \cdots \oplus R/\angles{p^{n_k(p)}}}$$
    with $1 \leq n_1(p) \leq \cdots \leq n_k(p)$, where all are uniquely determined. The prime powers $p^{n_j(p)}$ that occur are called the \textbf{elementary divisors} of $M$. (Note: the $\bigoplus$ above sums over some finite set of primes in $R$.)
\end{thm}

\begin{cor}[Fundamental Theorem of Finitely Generated Abelian Groups]
    If $G$ is any finitely generated abelian group, then
    $$G \cong \Z^m \oplus \Z/d_1\Z \oplus \cdots \oplus \Z / d_k\Z,$$
    with $d_1 \mid \cdots \mid d_k$. Moreover, $m,k,d_1,\dots,d_k$ are uniquely determined positive integers.
\end{cor}

\begin{ex}
    \begin{enumerate}
        \item Find all abelian groups of order 100: we have $100 = 2 \cdot 50 = 5 \cdot 20 = 10 \cdot 10$, giving $4$ groups in invariant factor form:
        $$\Z_{100}, \Z_2 \oplus \Z_{50}, \Z_5 \oplus \Z_{20},\Z _{10}\oplus \Z_{10}$$
        and in elementary divisor form:
        $$\Z_{2^2}\oplus \Z_{5^2},\Z_2 \oplus \Z_2 \oplus \Z_{5^2},\Z_5 \oplus \Z_5 \oplus \Z_{2^2}, \Z_2 \oplus \Z_2 \oplus \Z_5 \oplus \Z_5.$$
        \item Write $\Z_2 \oplus \Z_4 \oplus \Z_3 \oplus \Z_5 \oplus \Z_{25}$ in invariant form:
        $$\Z_{10}\oplus \Z_{300}$$
        \item For any fixed prime $p$, and $n \geq 1$, the number of isomorphism classes of abelian groups of order $p^n$ is $\pi (n)$, the number of partitions of $n$. For example, for $n=4$, we have $\pi (4)=5$, given by $4=3+1=2+2=2+1+1=1+1+1+1$.
    \end{enumerate}
\end{ex}

\begin{defn}
    An $R$-module $M$ is \textbf{cyclic} if it is generated by a single element $m \in M$. Equivalently, there exists a surjective linear map $\varphi:R \to M$, mapping $r \mapsto rm$. If $R$ is a PID, we deduce that $M \cong R / \ker \varphi \cong R/\angles a$ for some $a \in R$. So, the theorem we proved shows that a finitely generated module $M$ over a PID $R$ is a direct sum of cyclic modules.
\end{defn}

\section{Canonical Forms for Linear Operators}

Let $F$ be a field, $V$ an $F$-vector space, and $T:V  \to V$ an $F$-linear operator on $V$. Then, $V$ becomes an $F[x]$-module, by the action $f(x) \cdot v := f(T)(v)$ for any $f \in F[x], v \in V$.

Conversely, if $M$ is any $F[x]$-module, then $M$ is an $F$-module, i.e. an $F$-vector space. Also, the action of $X : M \to M$, $m \mapsto x \cdot m$ is an $F$-linear operator on $M$.

\begin{rmk}
    An $F[x]$-submodule of $M$ is just a $T$-invariant subspace of $M$.
\end{rmk}

\begin{ex}
    \begin{enumerate}
        \item Consider $V= F[x]$, the free $F[x]$ module of rank $1$. Then, $V$ is infinite dimensional as an $F$-vector space: $V \cong \bigoplus_{i=0}^\infty Fx^i$. The linear map $x : V \to V$ is the \textbf{shift operator}.
        \item Suppose $\dim _FV < \infty$, and $T \in \Hom (V,V)$, inducing an $F[x]$-module structure on $V$. Since $V$ is finite dimensional, it is a finitely generated $F[x]$-module, and $F[x]$ is a PID, so the structure theorem applies to $V$. Since $F[x]^k$ is infinite dimensional, $V$ must be a torsion $F[x]$-module. (This also follows from the fact that there exists a non-zero polynomial $p \in F[x]$ such that $p(T)=0$. Therefore, $p(x) \cdot v =0$ for all $v \in V$.)

        Consider $\Ann(V) \subset F[x]$. We must have $\Ann(V)  = \angles{m(x)}$ for a unique monic polynomial $m(x) \in F[x]$. In particular, $m(x)$ is the monic polynomial of least degree such that $m(T)=0$. Also, if $q(T)=0$ for some $q(x) \in F[x]$, then $m(x) \mid q(x)$. We call $m(x)$ the \textbf{minimal polynomial} of $T$.
    \end{enumerate}
\end{ex}

\begin{defn}
    Let $f(x) = x^n + a_{n-1}x^{n-1}+\cdots + a_1x+a_0$ be a monic polynomial. The \textbf{companion matrix} of $f$ is the $n \times n$ matrix
    $$C(f) = \matr{- a_{n-1} & 1 & 0  & \cdots & 0 \\ -a_{n-2} & 0 & 1 & \cdots & 0 \\ \vdots & \vdots & \ddots & \ddots & \vdots \\ \vdots & \vdots & \ddots & \ddots & 1 \\ -a_0 & 0 &\cdots & 0 & 0}$$
\end{defn}

\begin{proof}
    \begin{align*}
        p_{C(f)}(x) &= \det (xI_n - C(f)) = \det \matr{x + a_{n-1} & -1 \\ a_{n-2} & x-1 \\ \vdots & & \ddots \\ &&& -1 \\ a_0 & & & x} = \cdots \\ &= a_0 + x(a_1 + x(a_2 + \cdots + (a_{n-1}+x)\cdots ).
    \end{align*}
\end{proof}

\begin{rmk}
    Suppose $V$ is a finite dimensional $F$-vector space, and $T \in \Hom(V,V)$ giving an $F[x]$-module structure to $V$. Then, the structure theorem implies $V = W_1 \oplus \cdots \oplus W_K$ splits as a direct sum of $T$-invariant subspaces, with $W_i \cong F[x]/\angles{d_i(x)}$, where $d_i(x) \in F[x]$ are unique monic polynomials in $F[x]$ such that $d_1 \mid d_2 \mid \cdots \mid d_k$.
\end{rmk}

\begin{cor}
    \begin{enumerate}[(i)]
        \item  $V$ has an ordered basis $\mc C$ such that $(T)_\mc C$ has block diagonal form, where the blocks have size $\dim W_1,\dots, \dim W_K$
        \item The action of $T$ on $W_i$ corresponds to multiplication by $x$ in $F[x] / \angles {d_i}$. This allows us to choose a basis of $W_i$ so that $T\mid _{W_i}$ is represented by a simple matrix.
        \item $\Ann(V) = \angles {d_k}$. This implies the next proposition.
    \end{enumerate}
\end{cor}

\begin{prop}
    The minimal polynomial $m_T$ of $T$ equals the largest invariant factor of $V$. Moreover, all the invariant factors divide $m_T(x)$.
\end{prop}

Consider a single invariant subspace $W \cong F[x] / \angles d$, with
$$d(x) = x^r  + a_{r-1}x^{r-1}+\cdots  +a_1x + a_0,$$
and $a_i \in F$. Choose an ordered basis
$$(e_1,\dots,e_r) := ( \bar x^{r-1},\bar x^{r-2},\dots,\bar x,\bar 1)$$
of $F[x] / \angles d$. Let $\varphi : F[x] /\angles d \to W$ be the respective isomorphism. Then, the operator $X$ on $F[x]/ \angles d$ has the form:
\begin{align*}
    &xe_1 = \bar x ^r = -a_{r-1}e_1 - \cdots - a_0e_r \\ &xe_2 = e_1 \\ &\vdots \\ &x e_r = e_{r-1}.
\end{align*}
If $w_i = \varphi (e_i) \in W$ for $1 \leq i \leq r$, then $W$ has an ordered basis $\mathcal B := (w_1,\dots,w_r)$ with $T(w_1) = -a_{r-1}w_1 - \cdots - a_0w_r$, and $T(w_i) = w_{i-1}$ for $i \geq 2$. Hence we have
$$(T)_{\mc B} = \matr{-a_{r-1} & 1 \\ -a_{r-2} & 0 & 1 \\ \vdots & & 0 & \ddots \\ &&&\ddots&1 \\ -a_0 &&&& 0}=C(d).$$

The lemma implies that the characteristic polynomial of $T|_W =$ char polyn $(C(d)) = d(x)$. Applying this analysis to all the invariant subspaces $W_i$ of $T$, we obtain a basis $\mc C$ of $V$ such that
\begin{equation} \label{eq:ratcon}
    (T)_{\mc C} = \matr{C(d_1)  \\  & C(d_2) & 0 \\ &0& \ddots \\ &&& C(d_k)}
\end{equation}

\begin{defn}
    A matrix in the form \eqref{eq:ratcon} is said to be in \textbf{rational canonical form}. The polynomials $d_1\mid \cdots\mid d_k$ are the invariant factors of the matrix.
\end{defn}

\begin{thm}
    \begin{enumerate}[(i)]
        \item There is a basis $\mc C$ for $V$ with respect to which $(T)_{\mathcal C}$ is in rational canonical form.
        \item The rational canonical form for $T$ is unique. In particular, two linear transformations $V \to V$ are similar if and only if they have the same rational canonical form.
    \end{enumerate}
\end{thm}

\begin{proof}
    (Proving ii) Note that the process of construction the rational canonical form (R.C.F) of $T$ is invertible: if $\mathcal C$ is a basis of $V$ so that $(T)_{\mathcal C}$ is in rational canonical form, then $V \cong \bigoplus_{i=1}^k W_i$, and it is easy to check that each $W_i$ is a cyclic $F[x]$-module, with annihilator $d_{i}(x)$. Uniqueness now follows from the corresponding uniqueness statement in the structure theorem.
\end{proof}

\begin{cor}
    Let $A \in M_n(F)$ be an  $n \times n$ matrix.
    \begin{enumerate}[(i)]
        \item The characteristic polynomial of $A$ is the product $d_1\cdots d_k$ of all invariant factors of $A$.
        \item (Cayley-Hamilton Theorem) $m_A(x)$ divides the characteristic polynomial $p_A(x)$.
        \item The characteristic polynomial $p_A(x)$ divides some power of $m_A(x)$ (namely, $m_A^k$). In particular, these two polynomials have the same roots, not counting multiplicities.
    \end{enumerate}
\end{cor}

\begin{cor}
    Let $A \in M_n(F)$ be a nilpotent matrix. Then, $A$ is similar to an upper triangular matrix, whose non-zero entries are $1$'s adjacent to the main diagonal.
\end{cor}

\section{Jordan Canonical Form}

Assume that $F$ is algebraically closed, e.g. $F=\C$. Then, the elementary divisors of $V$ over $F[x]$ have the form $(x-\lambda)^r$ for some $\lambda \in F,r \geq 1$, while $V \cong \bigoplus _{i}F[x]/(x-\lambda_i)^{n_i}$ as an $F[x]$-module.

Suppose now that $W \cong F[x]/(x-\lambda)^r$ is a $T$-invariant subspace. Then, we choose a basis $(e_1,\dots,e_r)$ for this space with
$$e_1:=(\bar x - \lambda)^{r-1},\dots, e_{r-1}:=(\bar x - \lambda), e_r := 1.$$
Then, \begin{align*}
   & \bar x e_1 = \lambda e_1 + (\bar x - \lambda ) ^r = \lambda e_1 \\
   & \bar x e_2 = e_1 + \lambda e_2 \\ &\bar x e_3 = e_2 + \lambda e_3 \\ & \vdots \\ & \bar x e_r = e_{r-1}+ \lambda e_r.
\end{align*}
So, the action of $T$ on $W$ is given by the matrix
$$J_r(\lambda) := \matr{\lambda & 1 \\ & \lambda & 1 & 0 \\ && \ddots  & \ddots \\ &0&& \ddots & 1 \\ &&&& \lambda}$$

\begin{defn}
    $J_r(\lambda)$ is called a \textbf{Jordan block} of size $r$ with eigenvalue $\lambda$. A matrix is said to be in \textbf{Jordan canonical form} if it is a block diagonal matrix with Jordan blocks along the main diagonal.
\end{defn}

\begin{thm}
    Let $T \in \Hom(V,V)$, and assume $F$ contains all eigenvalues of $T$. Then, there exists a basis $\mathcal C$ of $V$ such that $(T)_{\mc C}$ is in Jordan canonical form. Moreover, the Jordan canonical form is unique, up to a permutation of the Jordan blocks along the diagonal.
\end{thm}

\section{Converting a Matrix Into Jordan Canonical Form}

\begin{rmk}
    Every Jordan block $J_k(\lambda)$ has a unique eigenvector (up to scale), namely $e_1 = (1,0,\dots,0)$, as $Te_1 = \lambda e_1$. Indeed, the $\lambda$-eigenspace $W_\lambda = \ker (T|_W - \lambda I) = \angles {e_1}$ is $1$-dimensional. This allows us to equate the dimension of the eigenspace $V_\lambda=\ker (T-\lambda I)$ corresponding to $\lambda$ with the number of Jordan blocks with eigenvalue $\lambda$.
\end{rmk}

\textbf{Case 1: The roots of the characteristic polynomial $p_T(x)$ of $T$ are distinct.}

Then, the Jordan form is diagonal.

\textbf{Case 2: The root $\alpha$ of $p_T(x)$ has multiplicity $r$.}

Here, there are various possibilities for the part of the Jordan matrix with $\alpha$ on the diagonal.

\begin{enumerate}
    \item[$r=2$:] $\matr{\alpha & 1 \\ & \alpha}$, $\matr{\alpha  \\ & \alpha}$
    \item [$r=3$:] $\matr{\alpha & 1 \\ & \alpha & 1 \\ && \alpha}$, $\matr{\alpha & 1 \\ & \alpha \\ && \alpha}$, $\matr{\alpha \\ & \alpha \\ && \alpha}$
    \item [$r=4$:] $\matr {\alpha & 1 \\ & \alpha & 1 \\ && \alpha & 1 \\ &&& \alpha & 1}, \matr{\alpha & 1 \\ & \alpha & 1 \\ && a \\ &&& \alpha}, \matr {\alpha & 1 \\ & \alpha \\ && \alpha & 1 \\ &&& \alpha}, \matr{\alpha & 1 \\ & \alpha \\ && \alpha\\ &&& \alpha},\matr{\alpha &  \\ & \alpha \\ && \alpha\\ &&& \alpha}$
\end{enumerate}

The cases $r=2$ and $r=3$ are distinguished by computing $\dim V_{\alpha}$. When $r=4$, we see that $\dim V_\alpha = \dim \ker (A - \alpha I)$ is given by $1,2,2,3,4$ respectively. The remaining $2$ cases are distinguished by computing $\dim (A -\alpha I)^2$ (the matrix $(A-\alpha I)^2$ is non-zero, or zero, respectively).

In general, one can show that the dimensions of the null spaces of the matrices $(A - \alpha I)^k$ for $k = 1,2,\dots, \frac{r}{2}$ distinguish the Jordan forms in all cases. The spaces $\ker (A - \alpha I)^k$ are called \textbf{generalized eigenspaces.}

\begin{rmk}[][label=rmk:cyclicrmk]
    If $M \cong R/\angles{a}$ is cyclic, for $R$ a PID, then $\angles a = \Ann(M)$.

    Moreover, any submodule $N \subset M$ is cyclic, with $\Ann(N) = \angles b$, where $b \mid a$. Conversely, for any divisor $b \mid a$, $M$ has a \textit{unique} submodule $N$, with $\Ann (N) = \angles b$
\end{rmk}

\begin{proof}
    Submodules of $M$ are given by ideals of $R/\angles a$, so they're of the form $J/I$, where $I \subset J \subset R$. Therefore, they're of the form $\angles{b'}/\angles {a}$, where $\angles a \subset\angles{b'}$, i.e. $b' \mid a$. If we then write $a = bb'$, then the map
    \begin{align*}
        Rb'/Ra &\to R/Rb=R/\angles b
        \\ rb' + Rbb' &\mapsto r+Rb
    \end{align*}
    is an $R$-module isomorphism, so $\Ann (N) = \angles b$, and the converse also follows.
\end{proof}

\section{Normal Linear Operators, and The Spectral Theorem}

Let $V$ be a finite dimensional complex vector space, equipped with a hermitian inner product $\angles ,$. Last semester (in the homework), we showed that for all $T \in \Hom(V,V)$, there exists a unique $T^* \in \Hom(V,V)$, called the \textbf{adjoint} of $T$, characterized by the property
$$\angles{Tx,y}=\angles{x,T^*y}\:\:\: \forall x,y\in V.$$
Moreover, if $A \in M_n(\C)$, the \textbf{adjoint matrix} is $A^* = \bar A^t$.

\begin{rmk}
    Our constructions/results have analogues for real inner product spaces $V$, with almost the same proofs.
\end{rmk}

\begin{lem}
    Let $T \in \Hom(V,V)$ and $\mc B = (b_1,\dots,b_n)$ be an ordered orthonormal basis of $V$. Moreover, suppose $A:= (T)_{\mc B}=\{a_{ij}\}$. Then, $(T^*)_B = A^*$.
\end{lem}

\begin{proof}
    Since $T(b_j) = \sum_{i=1}^na_{ij}b_i$, and $\mc B$ is orthonormal, we have
    \begin{align*}
        a_{ij}=\angles{T(b_j),b_i}=\angles{b_j,T^*(b_i)}=\overline{\angles{T^*b_i,b_j}}
        \\ \implies \angles{T^*(b_j),b_i}=\overline{a_{ji}}.
    \end{align*}
\end{proof}

\begin{rmk}
    The operation $^*$ on $M_n(\C)$ or on $\Hom(V,V)$ is an analogue of complex conjugation on $\C$. One can check that
    \begin{enumerate}[(i)]
        \item $(S+T)^*=S^* + T^*$
        \item $(\lambda T)^* = \bar \lambda T^*$
        \item $(ST)^*=T^*S^*$
        \item $(T^*)^*=T$
    \end{enumerate}
\end{rmk}

\begin{notebox*}[Motivating Question]
    Given $T \in \Hom(V,V)$, when is there an orthonormal basis of eigenvectors for $T$?
\end{notebox*}

\textbf{Necessary condition}: If $\mc B$ is orthonormal, and
$$(T)_{\mc B}=D= \matr{c_1 \\ & \ddots \\ && c_n}$$
is diagonal, then
$$(T^*)_{\mc B} = \bar D=\matr{\bar c_1 \\ & \ddots \\ && \bar c_n}.$$
If $V$ is a real inner product space, we have $c_i \in \R$, hence $T = T^*$. If $V$ is a $\C$-inner product space, we may have $T \neq T^*$, however $TT^*=T^*T.$

\begin{defn}
    $T \in \Hom(V,V)$ is \textbf{self-adjoint} if $T=T^*$, or \textbf{normal} if $TT^*=T^*T$
\end{defn}

\begin{prop}[][label=prop:sixprops]
    Let $V$ be a hermitian inner product space, $\dim_\C V < \infty$, and $T:V \to V$ a normal linear operator on $V$.
    \begin{enumerate}
        \item If $f \in \C[x]$ is any polynomial, then $f(T)$ is normal.
        \item $||Tv||=||T^*v||$ for all $v \in V$.
        \item $\ker T = (\Img T)^\perp$.
        \item $T^2v =0 \implies Tv=0$ for all $v \in V$.
        \item $v \in V$ is an eigenvector for $T$ with eigenvalue $\lambda \in \C \iff v$ is an eigenvector for $T^*$ with eigenvalue $\bar \lambda \in \C$.
        \item The eigenspaces of distinct eigenvalues of $T$ are orthogonal.
    \end{enumerate}
\end{prop}

\begin{proof}
    \hfill
    \begin{enumerate}[(i)]
        \item Since $(\lambda T^n)^* = \bar \lambda (T^*)^n$, this claim follows directly from the definition of normality.
        \item
        $$||Tv||^2 = \angles{Tv,Tv}=\angles{v,T^*Tv}= \angles{v,TT^*v}=\angles{T^*v,T^*v}=||T^*v||^2$$
        \item  \begin{align*}
            u \in (\Img T)^\perp \iff \angles{u, Tv}=0 \:\: \forall v\in V \iff \angles{T^*u,v}=0 \\ \iff T^*u=0 \overset{(ii)}{\iff} Tu=0 \iff u \in \ker T
        \end{align*}
        \item
        $$T^2v = 0 \implies Tv \in \Img T \cap \ker T \overset{(iii)}{=}\{0\}$$
        \item By (i), $T-\lambda I$ is normal for all $\lambda \in \C$. Then, by (ii),
        $$||(T-\lambda I)(v)|| = 0 \iff ||(T^*-\bar \lambda I)(v)||=0,$$
        hence $v$ is an eigenvector for $T$ with eigenvalue $\lambda \in \C \iff v$ is an eigenvector for $T^*$ with eigenvalue $\bar \lambda \in \C$.
        \item Suppose $Tv = \lambda v$ and $Tv' = \mu v'$, where $v,v' \neq 0$, and $\lambda \neq \mu$. Then, we compute
        \begin{align*}
            \angles{v,v'}=\angles{\frac{1}{\lambda - \mu}(T-\mu I)(v),v'}=\frac{1}{\lambda - \mu}\angles{v,(T-\mu I)^*(v')}=\frac{1}{\lambda - \mu}\angles{v,0}=0
        \end{align*}
        because $v'$ is an eigenvector for $T^*$ with eigenvalue $\bar \mu$, by (v).
    \end{enumerate}
\end{proof}

\begin{thm}[Spectral Theorem (Complex Case)]
    Suppose $V$ is a finite dimensional Hermitian space, and $T\in \Hom(V,V)$ is normal. Then, $V$ has an orthonormal basis of eigenvectors for $T$.
\end{thm}

\begin{proof}
    We claim that the minimal polynomial $m_T(x)$ is a product of distinct irreducible factors. Indeed, if $d(x)$ is an irreducible factor of $m_T(x)$, then we show that $d^2(x)$ does not divide $m_T(x)$.

    Suppose in fact that $d^2(x) \mid m_T(x)$. Recall that, as a $\C[x]$-module,
    $$V \cong \C[x]/\angles{m_T} \oplus \cdots.$$
    By \Cref{rmk:cyclicrmk}, we therefore have that $\C[x] / \angles{d^2}$ is a cyclic submodule of $V$. Hence, there is a vector $v \in V$ with $d^2(T)(v) =0 $, but $d(T)(v) \neq 0$. However by (i) from \Cref{prop:sixprops}, we know $d(T)$ is normal, but this contradicts (iv). Thus, as every irreducible polynomial in $\C[x]$ is linear, it follows that $m_T(x)$ is a product of distinct linear factors. Hence, $T$ is diagonalize (see HW), and the eigenspaces for $T$ are orthogonal by (vi). So, we choose an orthonormal basis for each eigenspace.
\end{proof}

\begin{thm}[Spectral Theorem (Real Case)]
    Suppose $V$ is a finite dimensional real inner product space, and $T:V \to V$ is a self-adjoint operator. Then, $V$ has an orthonormal basis of eigenvectors for $T$.
\end{thm}

\begin{proof}
    Let $V_\C := V \oplus V$, and we can make $V_\C$ into a $\C$-vector space by setting
    $$(a+bi)\cdot (u,v)=(au-bv,bu+av)$$
    for all $a,b \in \R$ and $u,v \in V$. Then, we define an inner product on $V_\C$ by
    $$\angles{(u,v),(u',v')}=\angles{u,u'}+\angles{v,v'}+i(\angles{v,u'}-\angles{u,v'})$$
    (one can check this). Next, extend $T$ to a linear map $T_\C :V_\C \to V_\C$ by $T_\C (u,v) := (T(u),T(v))$. One can check that $T_\C$ is $\C$-linear on $V_\C$, and also self-adjoint. By (v) from \Cref{prop:sixprops}, all the eigenvalues of $T_\C$ are real, and $m_{T_\C}(x)$ is a product of distinct (real) linear factors, hence $m_{T_\C} \in \R[x]$. If $f \in \R[x]$, then $f(T_\C)(u,v)=(f(T)(u),f(T)(v))$. Hence, $m_T(x) = m_{T_\C}(x)$, and we conclude that $T$ is diagonalizable. The rest of the proof is the same as in the complex case.
\end{proof}

\begin{rmk}
    We have $V_\C = \C \otimes _\R V$, and $T_\C = 1_\C \otimes _\R T$
\end{rmk}
\begin{proof}
    Indeed, note that
    $$C \otimes _{\R}V = (R \oplus i\R)\otimes_\R V \cong (1 \otimes V)\oplus (i \otimes V) = V \oplus iV,$$
    and
    $$(a+bi)(1 \otimes u + i \otimes v) = (a+bi) \otimes u + (ai-b)\otimes v = 1 \otimes (au-bv)+i\otimes (bu+av)$$
\end{proof}

\begin{thm}[Spectral Theorem Matrix Form]
    \begin{enumerate}[(a)]
        \item If $A \in M_n(\C)$ is a normal matrix, then there exists a unitary matrix $U$ such that $A=UDU^{-1}=UDU^*$, where $D$ is diagonal.
        \item If $A \in M_n(\R)$ is symmetric, then there exists an orthogonal matrix $P$ such that $A=PDP^{-1}=PDP^t$, where $D$ is diagonal.
    \end{enumerate}
\end{thm}

\begin{ex}
    $A=\matr{1 & 2  \\ 2 & 1}$ has eigenvalues $3$ and $-1$, with corresponding eigenvectors $(1,1)$ and $(-1,1)$, respectively. We deduce that
    $$A = P\matr{3 & 0 \\ 0 & -1}P^t,$$
    where
    $$P = \matr{\sqrt 2/2 & - \sqrt 2/2 \\ \sqrt 2 /2 & \sqrt 2/2}$$
\end{ex}

\section{Bilinear Forms}

\begin{defn}
    Let $V$ be a finite dimensional vector space over a field $F$, with basis $e_1,\dots,e_n$. Write
    $$x = \sum_{i=1}^nx_ie_i,y=\sum_{j=1}^ny_je_j$$
    with $x_i,y_j \in F$. A \textbf{bilinear form} on $V$ is a bilinear map $B:V \times V \to F$, such that
    $$B(x,y) = \sum_{i,j=1}^n x_iy_jB(e_i,e_j).$$
    The matrix $A:= \{ B(e_i,e_j)\}_{1\leq i , j \leq n}$ is called the matrix of $B$ relative to the basis $E = (e_1,\dots,e_n)$.

    Conversely, if $A= \{a_{ij}\}$ is any $n \times n$ matrix with entries in $F$, we define the function
    $$B_A(x,y):= x^tAy = \sum_{i,j=1}^na_{ij}x_iy_j.$$
    The determinant $\det \{B(e_i,e_j)\}$ is called the \textbf{discriminant} of $B$.
\end{defn}

Suppose we change the basis $E$ to another basis $E' = (e_1',\dots,e_n')$, and let $P:= (\id )_{E ,E'}$ be the change of basis matrix, so that if $P = \{p_{ij}\}$, then $e_i' = \sum_{j=1}^n p_{ij}e_j$ for all $i \in \{1,\dots,n\}$. Then,
\begin{align*}
    B(e_i',e_j') = B\paren{\sum_kp_{ik}e_k,\sum_l
p_{jl}e_l}= \sum_{k,l} p_{ik}B(e_k,e_l) p_{jl}.
\end{align*}
Hence, if $A = \{B(e_i,e_j)\}$ and $A'=\{B(e_i',e_j')\}$, then,
$$A'=PAP^t.$$
We therefore see that $\det (A')= (\det P)^2 \cdot \det (A)$, so the discriminant $\det A$ of $B$ is changed to $(\det P)^2 \cdot \det A$ upon changing the basis. That is, the discriminant \underline{does} depend on choice of basis.

\begin{defn}
    We define $V^L:= \{ x \in V \mid B(x,v)=0$ for all $v \in V \}$. We call this the \textbf{left radical} of $B$. Similarly, we have the \textbf{right radical} $V^R:= \{ y \in V \mid B(v,y)=0$ for all $v \in V \}$
\end{defn}

\begin{thm}
    The following are equivalent:
    \begin{enumerate}[(i)]
        \item $V^L=0$
        \item $V^R=0$
        \item The matrix $A$ of $B$ relative to any basis is invertible.
    \end{enumerate}
\end{thm}

\begin{proof}
    Let $e_1,\dots,e_n$ be any basis of $V$, and $A=\{B(e_i,e_j)\}$. Then, clearly
    \begin{align*}
        x \in V^L \iff B(x,e_i) =0 \:\:\forall i=1,2\dots,n \\ x \in V^R \iff B(e_i,x)=0 \: \: \forall i = 1,2,\dots,n.
    \end{align*}
    So, if we write $x = \sum x_je_j$, then $V^L=0$ (resp. $V^R=0$) is equivalent to saying that the linear system $A^tx = 0$ (resp. $Ax=0$) has only the zero solution. These are equivalent to the condition $\det A \neq 0$.
\end{proof}

\begin{defn}
    A bilinear form $B$ is called \textbf{non-degenerate} if its matrix with respect to any basis is invertible.
\end{defn}

Fix $y \in V$. Then, $y$ gives rise to a linear functional on $V$, namely $x \mapsto B(x,y)$. This functional only depends on the equivalence class of $y$ modulo $V^R$. Hence, we get a homomorphism $V \to V^*$ whose kernel is $V^R$ by definition. Hence, we obtain an injection $V/V^R \hookrightarrow V^*$. If $B$ is non-degenerate, we obtain an isomorphism $V \cong V^*$ by $y \mapsto (x \mapsto B(x,y))$. A similar analysis applies to the second variable and $V^L$.

\begin{prop}
    Suppose $B$ is non-degenerate and $\dim _FV<\infty$. Then, any linear functional on $V$ has the form $x \mapsto B(x,y)$ (and also $x \mapsto B(y,x)$) for some $y \in V$.
\end{prop}

\begin{defn}
    We say that $x$ is \textbf{orthogonal} to $y$ if $B(x,y)=0$. We denote this $x \perp y$.
\end{defn}

\begin{thm}
    The relation $\perp$ is symmetric $\iff$
    \begin{enumerate}
        \item $B$ is \textbf{alternating}: $B(x,x)=0$ OR
        \item $B$ is \textbf{symmetric}: $B(x,y)=B(y,x)$.
    \end{enumerate}
\end{thm}

\begin{proof}
    On homework.
\end{proof}

\section{Symmetric Bilinear Forms and Quadratic Forms}

Suppose $B:V \times V \to F$ is a symmetric bilinear form, and $\Char{F} \neq 2.$

\begin{defn}
    The corresponding \textbf{quadratic form} to $B$ is $Q(x) := B(x,x)$, which is a function $Q:V \to F$.
\end{defn}

\begin{rmk}
    We can recover $B$ from $Q$ by
    $$B(x,y) = \frac{1}{2}(Q(x+y)-Q(x)-Q(y))$$
\end{rmk}

By choosing a basis $E$ of $V$, we can represent $B$ by a symmetric matrix $A$ so that $B(x,y) = x^tAy$, and $Q(x) = x^tAx$.

\begin{ex}
    \begin{align*}
        Q(x,y,z) = x^2 + 2y^2 - z^2 + 10xy + 6xz - 4yz = \matr{x & y & z} \matr{1 & 5 & 3 \\ 5 & 2 & -2 \\ 3 & -2 & -1}\matr{x \\ y \\ z}
    \end{align*}
\end{ex}

\begin{defn}
    Suppose $Q=Q_A(x) = x^tAx$ for $A \in M_n(\R)$ symmetric. We say that $Q$ is
    \begin{enumerate}[(i)]
        \item \textbf{positive definite} if $Q(x)>0$ for all $x \neq 0$.
        \item \textbf{negative definite} if $Q(x)<0$ for all $x \neq 0$.
        \item \textbf{indefinite} if $Q(x)>0$ and $Q(x')<0$ for some $x,x'$.
    \end{enumerate}
\end{defn}

To determine which property $Q$ satisfies, we can use the spectral theorem.

Write $A= PDP^t$, where $P$ is orthogonal and $D=\matr{\lambda _ 1 \\ & \ddots \\ && \lambda _n}$ is diagonal. Then, we introduce new coordinates $y:=(x)_P=P^{-1}x \iff x=Py$. So,
\begin{align*}
    Q = x^tAx = (Py)^t A (Py) = y^t Dy
\end{align*}

\begin{ex}
    Let $A = \matr{1 & 2 \\ 2 & 1}$, so $Q_A(x_1,x_2) = x_1^2 + x_2^2 + 4x_1x_2$. We have $A=PDP^t$ with
    $$D = \matr{3 & 0 \\ 0 & -1}, P= \frac{\sqrt 2}{2}\matr{1 & -1 \\ 1 &1}.$$
    Then,
    $$y = \matr{y_1 \\ y_2}= P^t \matr{x_1 \\ x_2} = \frac{\sqrt 2}{2}\matr{x_1 + x_2 \\ x_1-x_2},$$
    so
    \begin{align*}
        Q(y) &= \matr{y_1 & y_2}\matr{3 & 0 \\ 0 & -1}\matr{y_1 \\ y_2}=3 \paren{\frac{\sqrt 2}{2}(x_1+x_2)}^2 - \paren{\frac{\sqrt 2}{2}(x_1-x_2)^2}
        \\ &= \frac{1}2{(3(x_1+x_2)^2 -(x_1-x_2)^2)}.
    \end{align*}
    Hence, $Q$ is indefinite.
\end{ex}

\begin{thm}
    For any real quadratic form $Q = x^tAx$ in $n$ variables, there is an orthogonal change of basis in which that form becomes
    $$Q(y_1,\dots,y_n) = \lambda _1y_1^2 + \cdots + \lambda _ny_n^2,\:\: \lambda _i \in \R.$$
    $Q$ is positive (resp. negative) definite if $\lambda _i >0$ (resp. $\lambda _i < 0$) for all $i=1,\dots,n$. $Q$ is indefinite if $\lambda _i >0$ and $\lambda _j<0$ for some $i,j$.
\end{thm}

\begin{defn}
    We say that $Q$ is \textbf{positive semi-definite} (resp. \textbf{negative semi-definite}) if $Q(x) \geq 0$ (resp. $Q(x) \leq 0$) for all $x \in \R^n$.
\end{defn}

Over any field $F$, one can show:
\begin{thm}
    Let $B$ be a symmetric bilinear form on a vector space $V$ over a field $F$ with $\Char F \neq 2$. Then, $V$ has a basis $E = (e_1,\dots,e_n)$ such that the matrix of $B$ relative to $E$ has the form
    $$\matr{\lambda _ 1 \\ & \ddots \\ && \lambda _r \\ &&&0 \\ &&&&\ddots\\ &&&&&0}$$
    where $\lambda _i \in F \setminus\{0\}$ for all $i\in \{1,\dots,r\}$. Note that $B(e_i,e_j)=0$ for all $i \neq j$, so $E$ is an \textbf{orthogonal} basis of $V$.
\end{thm}

\section{Alternating Bilinear Forms}

Let $V$ be a vector space over $F$, and $B:V \times V \to F$ an alternating ($B(x,x)=0$) bilinear form on $V$. Then, $B(x,y) = -B(y,x)$ for all $x,y \in V$.

If $E=(e_1,\dots,e_n)$ is any basis of $V$, we see that $B(e_i,e_j) = - B(e_j,e_i)$ for $i \neq j$, and $B(e_i,e_i)=0$ for all $i$. Hence, if $A = \{ B(e_i,e_j)\}_{i,j}$, then $A + A^t=0$, i.e. $A$ is \textbf{skew-symmetric}.

Conversely, if $A$ is a skew-symmetric matrix in $M_n(F)$, then $B(x,y):= x^tAy$ is an alternating form, since
$$B(x,x) = \sum_{i,j} a_{ij}x_ix_j= \sum_{i <j}(a_{ij}+a_{ji})x_ix_j=0.$$

\begin{thm}
    If $B$ is an alternating bilinear form, then there exists a basis of $V$ such that the matrix of $B$ relative to this basis has the form
    $$A':= \matr{S \\ & \ddots \\ && S \\ &&& 0 \\ &&&& \ddots \\ &&&&& 0},$$
    where $S= \matr{0 & 1 \\ -1 & 0}$.
\end{thm}

\begin{proof}
    If $B=0$, the result is trivial, so assume $B \neq 0$. Then, we can find $u,v \in V$ such that $B(u,v) \neq 0$. Since $B(u, \lambda u)=0$ for all $\lambda \in F$, we know $u$ and $v$ are linearly independent. Say $B(u,v) = \lambda \neq 0$, and define $u_1:=u,v_1 := \lambda^{-1}v$. Then, $B(u_1,v_1) = 1 = -B(v_1,u_1)$.

    Inductively, suppose we found linearly independent vectors $E_k:= (u_1,v_1,\dots,u_k,v_k)$ such that $B(u_i, v_i)= 1 = -B(v_i,u_i)$ for all $i \in \{1,\dots,k\}$, and $B(z,w)=0$ for every other pair $z,w \in E_k$.

    Let $V_k:= \Span \{ u_1,v_1,\dots,u_k,v_k\}$. We claim $V = V_k \oplus V_k^\perp$. Indeed, since the matrix $B|_{V_k}$ is the diagonal matrix with $k$ $S$'s, which is invertible, we know that $B|_{V_k}$ is non-degenerate. So, $V_k \cap V_k^\perp =0$. Moreover, if $x \in V$, we may define $y := x- \sum_{i=1}^kB(x,v_i)u_i + \sum_{i=1}^kB(x,u_i)v_i$. Then, check that $B(y,u_j)=B(y,v_j)=0$ for all $j \in \{1,\dots,k\}$. Hence $y \in V_k^\perp$, so it follows that $x \in  V_k \oplus V_k^\perp$.

    Now, if $B|_{V_k^\perp}=0$, then choose a basis $w_1,\dots,w_{n-2k}$ for $V_k^\perp$, and append it to $E_k$ to obtain the required basis. Otherwise, we can choose $u_{k+1},v_{k+1} \in V_k^\perp$ such that $B(u_{k+1},v_{k+1})=1 = -B(v_{k+1},u_{k+1})$, and then proceed inductively.
\end{proof}

\begin{cor}
    The rank of a skew-symmetric matrix with entries in $F$ is even, and its determinant is a square.
\end{cor}

\begin{proof}
    If $A + A^t=0$, and $B_A(x,y)=x^tAy$, then the previous theorem gives $A=PA'P^t$, where $A'$ is the matrix from the theorem, and $P$ is a change of basis matrix. So, $P$ is invertible, hence $A$ and $A'$ have the same rank, so $A$ has even rank. Also, $\det (A) = (\det P)^2 \cdot \det (A')=(\det P)^2$ (or $0$).
\end{proof}

\section{The Pfaffian}

If $A = \{ a_{ij}\}_{1 \leq i < j \leq n}$ is skew-symmetric, then $\det A$ is a polynomial in the variables $a_{i,j}$ for $i<j$, i.e. in $R:= \Z[a_{ij},i<j]$.

\begin{ex}
    If $n$ is odd, then $\det A=0$. We can also see this as
    \begin{align*}
        \det A = \det (A^t) = \det (-A) = (-1)^n\det A
    \end{align*}
\end{ex}

\begin{ex}
    $$\det\matr{0 & a_{12} & a_{13} & a_{14} \\ & 0 & a_{23} & a_{24} \\ &&0  & a_{34} \\ &&&0}=(a_{12}a_{34}-a_{13}a_{24}+a_{14}a_{23})^2.$$
    We'll call the polynomial inside the square the Pfaffian.
\end{ex}

The quotient field of $R$ is $\Q (a_{ij},i<j)$. The corollary implies $\det A = \paren{\frac{f}{g}}^2$ for some $f,g \in R$. Therefore, $f^2 = \det (A) \cdot g^2$, hence $g^2 \mid f^2$ in $R$, so $g \mid f$ in $R$. That is, $\frac{f}{g}=h$ for some $h \in R$, hence $\det A = h^2$, and $h$ is determined up to a sign $\pm 1$ by this relation. When $A=A'=\w{diag}(S,S,\dots,S)$, we obtain $1= \det A'= h(a'_{12}=1,\dots)^2$. We fix the sign of $h$ so that $h(a_{12}'=1,\dots)=1$.

\begin{defn}
    The resulting polynomial from the discussion above is denoted $\Pf (A)$, the \textbf{Pfaffian} of $A$.
\end{defn}

In general,
\begin{align*}
    \Pf (A) = \sum \sgn (\sigma)a_{\sigma ( 1) \sigma (2)}\cdots a_{\sigma (2n-1)\sigma (2n)}
\end{align*}
summed over all $\sigma \in S_{2n}$ such that $\sigma (2r-1)<\sigma (2r)$ for all $r$, and $\sigma(1) < \sigma (3) < \cdots  < \sigma ( 2n-1)$. There are $(2n-1)!! = (2n-1)(2n-3)\cdots 3\cdot 1$ terms in this sum.

There also exists a Laplace type expansion
\begin{align*}
    \Pf A = a_{12}\Pf A_{12}-a_{13}\Pf A_{13}+\cdots + (-1)^n a_{1n}\Pf A_{1n}
\end{align*}
